\documentclass[11pt]{article}
\usepackage[letterpaper,margin=0.85in]{geometry}\usepackage[T1]{fontenc}\usepackage[utf8]{inputenc}\usepackage{amsmath,amssymb,enumitem,booktabs,array,xcolor,fancyhdr,tabularx}
\setlist{itemsep=2pt,topsep=3pt}\pagestyle{fancy}\fancyhf{}\renewcommand{\headrulewidth}{0.3pt}
\fancyhead[L]{\small 3.5 Higher-Order Derivatives}\fancyhead[R]{\small Make-up key}\fancyfoot[C]{\small\thepage}
\definecolor{ans}{RGB}{28,125,74}\newcommand{\A}[1]{\textcolor{ans}{\textbf{#1}}}
\begin{document}

\section*{3.5 Higher-Order Derivatives}
\textbf{You missed this lesson. Copy this into your notes. That is all you need to do.}\par\medskip
\subsection*{Question of the day}
What does it mean to take the derivative of a derivative?\par
\subsection*{Targets}\begin{itemize}\item I can name f?, f? and f? and say what each one tells me.
\item I can find a second or higher derivative with the usual rules.
\item I can find d\textsuperscript{2}y/dx\textsuperscript{2} at a point on an implicit curve.\end{itemize}
\subsection*{Concepts}\begin{itemize}
\item \textbf{Notation} (f???(x) = d?y/dx?): f?, f?, f? are the first, second and third derivatives. d\textsuperscript{2}y/dx\textsuperscript{2} is the same as f? (or y?).
\item \textbf{The second derivative} (f? = (f?)?): f? is the derivative of f?. It is the slope of f?: how the slope of f is changing.
\item \textbf{Reading the sign} (f?(x) = slope of f?): f? > 0 means f? is increasing. f? < 0 means f? is decreasing.
\item \textbf{Patterns save work} (d?/dx? e\textasciicircum{}(kx) = k? e\textasciicircum{}(kx)): A polynomial loses one degree each time. sin and cos repeat every 4 derivatives. Each derivative of e\textasciicircum{}(kx) multiplies by k.
\item \textbf{Implicit curves} (y? = d/dx (dy/dx)): Differentiate dy/dx again. Each y term brings a new dy/dx. Then put in the dy/dx you already found.
\end{itemize}
\subsection*{What the class predicted, and what is true}\begin{itemize}
\item The derivative of $x^2$ is $2x$. The derivative of $2x$ is... \A{2.} The line $2x$ has slope 2, so differentiating again gives 2.
\item Keep differentiating $\sin x$: $\sin x$, $\cos x$, $-\sin x$, $-\cos x$, and then... \A{sin x.} The next derivative of $-\cos x$ is $\sin x$. The pattern repeats every 4 derivatives.
\end{itemize}
\section*{Huddle notes (answers filled in)}
\subsection*{Definitions}\begin{itemize}
\item \textbf{Higher derivative:} Differentiate \A{again}. The second derivative is f?(x) or d\textsuperscript{2}y/dx\textsuperscript{2}. The third is f?(x).
\item \textbf{What f? means:} f? is the \A{slope} of f?. It tells how the slope of f is \A{changing}.
\item \textbf{Reading the sign:} If f? > 0, the slope of f is \A{increasing}. If f? < 0, it is \A{decreasing}.
\item \textbf{Patterns:} A polynomial loses \A{one} degree each time. sin x repeats every \A{4} derivatives.
\end{itemize}
\subsection*{Worked example}
Let f(x) = x? -- 3x\textsuperscript{2} + 2x. Find f? and f?.\\
\begin{itemize}
\item First derivative: f?(x) = 4x\textsuperscript{3} -- 6x + 2
\item Second derivative: f?(x) = \A{12}x\textsuperscript{2} -- 6
\item Third derivative: f?(x) = \A{24}x
\item The fourth derivative is the constant \A{24}
\end{itemize}
\subsection*{Label the three graphs}
\begin{enumerate}
\item The cubic is the graph of \A{f}.
\item The parabola is the graph of \A{f?}.
\item The straight line is the graph of \A{f?}.
\end{enumerate}
\section*{One worked sample of each skill}
\subsection*{Second derivative}
\noindent\textbf{Level 1.} Find the second derivative.\\ \A{Answer: $24 x + 8$} \quad $f'(x)=12x^{2}+8x+3$, so $f''(x)=24x+8$.\par\medskip
\noindent\textbf{Level 2.} Find the second derivative.\\ \A{Answer: $4 e^{2 x}$} \quad $f'(x)=2e^{2x}$. Differentiating that again gives the answer.\par\medskip
\noindent\textbf{Level 3.} Find the second derivative.\\ \A{Answer: $\left(2 + 4 x^{2}\right) e^{x^{2}}$} \quad $f'(x)=2x\,e^{x^{2}}$. Differentiating that again gives the answer.\par\medskip
\noindent\textbf{Level 4.} Find the second derivative.\\ \A{Answer: $\left(x^{2} + 4 x + 2\right) e^{x}$} \quad $f'(x)=(x^{2}+2x)\,e^{x}$. Differentiating that again gives the answer.\par\medskip
\subsection*{Higher derivative at a point}
\noindent\textbf{Level 1.} $f(x)=e^{2x}$. Find $f^{(4)}(0)$.\\ \A{Answer: 16} \quad Each derivative multiplies by 2. After 4 derivatives: $2^{4}e^{2x}$. At $x=0$ that is $16$.\par\medskip
\noindent\textbf{Level 2.} $f(x)=\sin x$. Find $f^{(99)}(0)$.\\ \A{Answer: -1} \quad The pattern repeats every 4. 99 leaves remainder 3, so $f^{(99)}$ is the same as $f^{(3)}$. At $x=0$ the value is $-1$.\par\medskip
\noindent\textbf{Level 3.} $f(x)=\cos(2x)$. Find $f^{(6)}(0)$.\\ \A{Answer: -64} \quad Each derivative multiplies by 2, so the size is $2^{6}=64$. The sign follows the cos pattern: $-$. The answer is $-64$.\par\medskip
\noindent\textbf{Level 4.} $f(x)=\cos(3x)$. Find $f^{(4)}(0)$.\\ \A{Answer: 81} \quad Each derivative multiplies by 3, so the size is $3^{4}=81$. The sign follows the cos pattern: $+$. The answer is $81$.\par\medskip
\subsection*{Read the second derivative}
\noindent\textbf{Level 1.} $f''(x)<0$ on an interval means the slope of $f$ is...\\ (A) increasing \quad \A{(B) decreasing} \quad (C) constant \quad (D) undefined\\ \A{Answer: B} \quad $f''$ is the slope of $f'$. Positive: $f'$ goes up. Negative: $f'$ goes down. Zero: $f'$ stays level.\par\medskip
\noindent\textbf{Level 2.} $f'(x)=x^{2}-4x-4$. At $x=4$, is $f''$ positive, negative or zero?\\ \A{(A) positive} \quad (B) negative \quad (C) zero \quad (D) cannot tell\\ \A{Answer: A} \quad $f''(x)=2x-4$. At $x=4$ it equals $4$, which is positive.\par\medskip
\noindent\textbf{Level 3.} The table gives two values of $f'$. Estimate $f''(4.5)$ with the slope between them.\\ \A{Answer: -5} \quad Rise over run: $\dfrac{-34-(-9)}{7-2}=\dfrac{-25}{5}=-5$.\par\medskip
\noindent\textbf{Level 4.} The table gives two values of $f'$. Estimate $f''(2.5)$ with the slope between them.\\ \A{Answer: 6} \quad Rise over run: $\dfrac{36-(6)}{5-0}=\dfrac{30}{5}=6$.\par\medskip
\subsection*{Second derivative of a curve}
\noindent\textbf{Level 1.} The curve $x^2+y^2=16$ passes through $(0,4)$. Find $\dfrac{d^2y}{dx^2}$ at that point.\\ \A{Answer: -0.25} \quad $y'=-\dfrac xy=0$ at the point. Then $y''=-\dfrac{y-xy'}{y^2}=-\dfrac{y}{y^2}=-\dfrac1y$ there, which is $-0.25$.\par\medskip
\noindent\textbf{Level 2.} The curve $xy=4$ passes through $(1,4)$. Find $\dfrac{d^2y}{dx^2}$ at that point.\\ \A{Answer: 8} \quad $y'=-\dfrac yx=-4$ at the point. Differentiating again: $y''=-\dfrac{y'x-y}{x^2}=8$.\par\medskip
\noindent\textbf{Level 3.} The curve $x^3+y^3=9$ passes through $(2,1)$. Find $\dfrac{d^2y}{dx^2}$ at that point.\\ \A{Answer: -36} \quad Find $y'$ first, differentiate it again (each $y$ brings a new $y'$), then put in $y'$ and the point. The answer is $-36$.\par\medskip
\noindent\textbf{Level 4.} The curve $x^3+y^3=9$ passes through $(2,1)$. Find $\dfrac{d^2y}{dx^2}$ at that point.\\ \A{Answer: -36} \quad Find $y'$ first, differentiate it again (each $y$ brings a new $y'$), then put in $y'$ and the point. The answer is $-36$.\par\medskip
\end{document}